\chapter{Spinor and Gamma-Matrix Conventions}
\label{app:spinor-conventions}

This appendix records the spinor conventions used throughout the monograph.
The convention is the Weinberg-compatible four-dimensional gamma-matrix
basis, adapted to the mostly-plus metric used in this monograph.  The purpose
of the appendix is not merely notational: signs in the Dirac adjoint, current
normalization, chirality, Majorana conjugation, anomaly traces, and
supersymmetry formulae are all fixed by the same data.  All comparisons below
keep the Lorentzian metric fixed as \(\eta=\operatorname{diag}(-,+,+,+)\).

\section{Four-Dimensional Lorentzian Clifford Algebra}

The Lorentzian metric is
\[
  \eta_{\mu\nu}=\operatorname{diag}(-,+,+,+),
  \qquad
  \eta^{\mu\rho}\eta_{\rho\nu}=\delta^\mu{}_\nu .
\]
Gamma matrices are endomorphisms of a complex four-dimensional vector space
\(S\) and obey
\begin{equation}
  \{\gamma^\mu,\gamma^\nu\}=2\eta^{\mu\nu}\mathbf 1_S .
  \label{eq:appendix-lorentz-clifford}
\end{equation}
Consequently
\[
  (\gamma^0)^2=-\mathbf 1_S,
  \qquad
  (\gamma^i)^2=\mathbf 1_S,\quad i=1,2,3 .
\]
For a covector \(p_\mu\) we write
\[
  \not p=p_\mu\gamma^\mu .
\]
If \(p^2=\eta^{\mu\nu}p_\mu p_\nu\), then
\[
  \not p^{\,2}=p^2\mathbf 1_S .
\]

The explicit basis used for sign checks in the monograph is
\begin{equation}
  \gamma^0
  =
  -\ii
  \begin{pmatrix}
    0&I_2\\ I_2&0
  \end{pmatrix},
  \qquad
  \gamma^j
  =
  -\ii
  \begin{pmatrix}
    0&\sigma_j\\ -\sigma_j&0
  \end{pmatrix},
  \qquad j=1,2,3,
  \label{eq:appendix-weinberg-gamma-basis}
\end{equation}
where
\[
  \sigma_1=\begin{pmatrix}0&1\\1&0\end{pmatrix},
  \quad
  \sigma_2=\begin{pmatrix}0&-\ii\\ \ii&0\end{pmatrix},
  \quad
  \sigma_3=\begin{pmatrix}1&0\\0&-1\end{pmatrix}.
\]
The hermiticity identities are
\[
  (\gamma^0)^\dagger=-\gamma^0,
  \qquad
  (\gamma^i)^\dagger=\gamma^i .
\]
Define
\begin{equation}
  \beta=\ii\gamma^0,
  \qquad
  \bar\psi=\psi^\dagger\beta .
  \label{eq:appendix-dirac-adjoint}
\end{equation}
Then
\[
  \beta^\dagger=\beta,
  \qquad
  \beta^2=\mathbf 1_S,
  \qquad
  (\gamma^\mu)^\dagger\beta=-\beta\gamma^\mu .
\]
Thus the bilinear \(\bar\psi\psi\) is Lorentz invariant and
\(\bar\psi\gamma^\mu\psi\) transforms as a Lorentz vector, with the Lorentz
action defined below.  In these conventions the free Lorentzian Dirac density
used in the monograph is
\[
  \mathcal L_{\rm D}=-\bar\psi(\gamma^\mu\partial_\mu+m)\psi .
\]
For a real Abelian charge \(q\), the Hermitian vector current is
\[
  j^\mu=\ii q\,\bar\psi\gamma^\mu\psi .
\]
The explicit factor of \(\ii\) is not optional: it follows from
\eqref{eq:appendix-dirac-adjoint} and the anti-Hermitian nature of
\(\bar\psi\gamma^\mu\psi\) in the mostly-plus convention.

\section{Spin Representation and Chirality}

The spinor Lorentz generators are
\begin{equation}
  S^{\mu\nu}
  =
  -\frac{\ii}{4}[\gamma^\mu,\gamma^\nu].
  \label{eq:appendix-spin-generators}
\end{equation}
For \(\Lambda=\exp\omega\in SO^+(1,3)\), with
\(\omega_{\mu\nu}=-\omega_{\nu\mu}\), the corresponding element in the
spinor representation is
\[
  R(\Lambda)
  =
  \exp\!\left(-\frac{\ii}{2}\omega_{\mu\nu}S^{\mu\nu}\right).
\]
The Clifford relation gives
\[
  [S^{\mu\nu},\gamma^\rho]
  =
  -\ii\gamma^\mu\eta^{\nu\rho}
  +\ii\gamma^\nu\eta^{\mu\rho},
\]
and hence
\[
  R(\Lambda)\gamma^\mu R(\Lambda)^{-1}
  =
  \Lambda^\mu{}_\nu\gamma^\nu .
\]
The \(\beta\)-pairing is invariant:
\[
  R(\Lambda)^\dagger\beta R(\Lambda)=\beta .
\]

The chirality matrix is
\begin{equation}
  \gamma_5=-\ii\gamma^0\gamma^1\gamma^2\gamma^3 .
  \label{eq:appendix-gamma-five}
\end{equation}
It obeys
\[
  \gamma_5^2=\mathbf 1_S,
  \qquad
  \{\gamma_5,\gamma^\mu\}=0,
  \qquad
  [\gamma_5,S^{\mu\nu}]=0 .
\]
In the explicit basis \eqref{eq:appendix-weinberg-gamma-basis},
\[
  \gamma_5=
  \begin{pmatrix}
    I_2&0\\0&-I_2
  \end{pmatrix},
  \qquad
  P_\pm=\frac{1\pm\gamma_5}{2}.
\]
Thus a spinor \(s\in S\) decomposes as
\[
  s=
  \begin{pmatrix}\chi_a\\ \widetilde\chi_{\dot a}\end{pmatrix},
  \qquad
  \chi\in S_+:=P_+S,
  \qquad
  \widetilde\chi\in S_-:=P_-S,
  \qquad
  a,\dot a=1,2 .
\]
The matrix \(\gamma^\mu\) maps \(S_+\) to \(S_-\) and \(S_-\) to \(S_+\).

With the orientation convention
\[
  \epsilon^{0123}=+1,
\]
the four-dimensional trace convention is
\begin{equation}
  \operatorname{tr}
  \left(\gamma_5\gamma^\mu\gamma^\nu\gamma^\rho\gamma^\sigma\right)
  =
  4\ii\,\epsilon^{\mu\nu\rho\sigma}.
  \label{eq:appendix-gamma-five-trace}
\end{equation}
This is the convention used in the anomaly chapter.

\section{Two-Component Blocks and Index Conventions}

Let
\[
  \gamma^\mu
  =
  \begin{pmatrix}
    0&\rho^\mu_{+-}\\
    \rho^\mu_{-+}&0
  \end{pmatrix}
\]
with respect to \(S=S_+\oplus S_-\).  In the basis
\eqref{eq:appendix-weinberg-gamma-basis},
\[
  \rho^0_{+-}=-\ii I_2,
  \qquad
  \rho^j_{+-}=-\ii\sigma_j,
  \qquad
  \rho^0_{-+}=-\ii I_2,
  \qquad
  \rho^j_{-+}=+\ii\sigma_j .
\]
Equivalently,
\[
  (\gamma^\mu s)_a=(\rho^\mu_{+-})_a{}^{\dot b}
  \widetilde\chi_{\dot b},
  \qquad
  (\gamma^\mu s)_{\dot a}=(\rho^\mu_{-+})_{\dot a}{}^b\chi_b .
\]

Spinor \(SU(2)\) indices are raised and lowered by
\[
  \chi^a=\epsilon^{ab}\chi_b,
  \qquad
  \widetilde\chi^{\dot a}
  =
  \epsilon^{\dot a\dot b}\widetilde\chi_{\dot b},
  \qquad
  \epsilon^{12}=\epsilon_{12}=1,
  \qquad
  \epsilon^{\dot1\dot2}=\epsilon_{\dot1\dot2}=1 .
\]
With this convention, the two chiral blocks obey the sign relation
\begin{equation}
  (\rho^\mu_{+-})_a{}^{\dot b}
  =
  -(\rho^\mu_{-+})^{\dot b}{}_a
  =
  -\epsilon^{\dot b\dot c}
  (\rho^\mu_{-+})_{\dot c}{}^d
  \epsilon_{da}.
  \label{eq:appendix-rho-sign-relation}
\end{equation}
In lowered form,
\[
  (\rho_\mu)_{a\dot b}(\rho_\nu)^{a\dot b}
  =
  2\eta_{\mu\nu}.
\]
Complex conjugation exchanges the chiral and antichiral components:
\[
  \overline\chi{}^{\dot a}=(\chi_a)^*,
  \qquad
  \overline{\widetilde\chi}{}^{a}
  =
  (\widetilde\chi_{\dot a})^* .
\]

\section{Majorana Conjugation}

A Majorana condition is a Lorentz-covariant reality condition on the complex
spinor representation.  In the explicit basis
\eqref{eq:appendix-weinberg-gamma-basis}, set
\begin{equation}
  B=\gamma_2 .
  \label{eq:appendix-majorana-B}
\end{equation}
Here \(\gamma_2=\eta_{2\nu}\gamma^\nu=\gamma^2\).  Directly from the Pauli
matrices,
\[
  (\gamma^\mu)^*=B\gamma^\mu B^{-1},
  \qquad
  B^*B=\mathbf 1_S .
\]
For a spinor \(s\), define
\[
  s^c=B^{-1}s^* .
\]
Then \(s^c\) transforms in the same Lorentz representation as \(s\), and
\((s^c)^c=s\).  A Majorana spinor is a spinor satisfying
\[
  s=s^c,
  \qquad\hbox{equivalently}\qquad
  s^*=Bs .
\]
If \(\chi=P_+s\) is the chiral component, then the Majorana spinor may be
written
\[
  s=\chi+B^*\chi^* .
\]
The charge-conjugation matrix in the convention often used for transposed
spinor contractions is
\[
  C=B\gamma^0 .
\]
For a Majorana spinor \(s\),
\[
  \bar s=s^\dagger\ii\gamma^0
  =
  \ii s^T C .
\]
For two Grassmann-even Majorana spinors \(s_1,s_2\), with chiral components
\(\chi_1,\chi_2\), the Lorentz vector bilinear decomposes as
\[
  \bar s_1\gamma^\mu s_2
  =
  \overline\chi_1\,\rho^\mu_{+-}\chi_2
  +
  \overline\chi_2\,\rho^\mu_{+-}\chi_1,
\]
where the precise matrix placement is the one fixed in
\eqref{eq:appendix-rho-sign-relation}.  For Grassmann-odd spinors the same
formula must be combined with the sign rules for moving odd coefficients
through one another.

\section{Comparison with Wess-Bagger-Type Four-Dimensional Conventions}

At fixed Lorentzian metric \(\eta=\operatorname{diag}(-,+,+,+)\), a common
chiral gamma-matrix basis, used for example in Wess-Bagger-style supersymmetry
formulae, is
\begin{equation}
  \gamma^0_{\rm WB}
  =
  \begin{pmatrix}
    0&I_2\\ -I_2&0
  \end{pmatrix},
  \qquad
  \gamma^j_{\rm WB}
  =
  \begin{pmatrix}
    0&\sigma_j\\ \sigma_j&0
  \end{pmatrix}.
  \label{eq:appendix-wb-gamma-basis}
\end{equation}
These matrices obey the same Clifford relation
\[
  \{\gamma^\mu_{\rm WB},\gamma^\nu_{\rm WB}\}=2\eta^{\mu\nu}.
\]
The Weinberg-compatible basis used in the monograph is related to
\eqref{eq:appendix-wb-gamma-basis} by the chiral phase change
\begin{equation}
  U_{\rm ch}
  =
  \begin{pmatrix}
    I_2&0\\0&\ii I_2
  \end{pmatrix},
  \qquad
  \gamma^\mu=U_{\rm ch}\gamma^\mu_{\rm WB}U_{\rm ch}^{-1}.
  \label{eq:appendix-wb-phase-change}
\end{equation}
Thus if \(s_{\rm WB}\) and \(s\) are coordinate columns for the same abstract
spinor in the two bases, then
\[
  s=U_{\rm ch}s_{\rm WB}.
\]
The chiral component is unchanged, while the antichiral component is
multiplied by \(\ii\).  This basis change is the source of many sign
differences in two-component supersymmetry formulae.  In particular, with the
index-raising convention used in this appendix, the monograph has the minus
sign in \eqref{eq:appendix-rho-sign-relation}; in the
Wess-Bagger-type basis one often writes the corresponding relation without
that minus sign, together with a different convention for upper and lower
\(\epsilon\)-symbols and dotted-index contractions.  Formulae copied between
the two systems must therefore translate both the gamma matrices and the
index algebra.

The chirality matrix defined by
\[
  \gamma_5=-\ii\gamma^0\gamma^1\gamma^2\gamma^3
\]
is unchanged by the similarity transformation
\eqref{eq:appendix-wb-phase-change}.  The Dirac adjoint is also the same
abstract Hermitian pairing after the basis change:
\[
  \bar s=s^\dagger(\ii\gamma^0)
  =
  s_{\rm WB}^\dagger
  \left(U_{\rm ch}^\dagger \ii\gamma^0 U_{\rm ch}\right)
  s_{\rm WB}.
\]
The displayed matrix in the last expression is the Wess-Bagger-basis matrix
for the same invariant pairing.

\section{Dimensional Continuation and \texorpdfstring{\(\gamma_5\)}{gamma5}}

Dimensional regularization and dimensional reduction are perturbative
schemes; they are not definitions of a non-perturbative fermionic path
integral.  Whenever a chiral trace is present, the finite-dimensional
four-dimensional data must be specified before continuing loop momenta.
The convention used in the anomaly chapter is:
\[
  \gamma_5=-\ii\gamma^0\gamma^1\gamma^2\gamma^3
\]
is a strictly four-dimensional matrix.  In a split
\(D=4-\varepsilon\) momentum algebra, write
\[
  \ell=\ell_\parallel+\ell_\perp,
\]
where \(\ell_\parallel\) lies in the physical four-dimensional subspace and
\(\ell_\perp\) lies in the continued complement.  Then
\[
  \{\gamma_5,\not\ell_\parallel\}=0,
  \qquad
  [\gamma_5,\not\ell_\perp]=0 .
\]
Together with \eqref{eq:appendix-gamma-five-trace}, these rules define the
chiral trace algebra used in the dimensional-reduction anomaly calculation.
Changing this prescription changes intermediate contact terms and must be
compensated by an explicit finite local counterterm if the same renormalized
Ward identity is to be recovered.

\section{Three-Dimensional Lorentzian Spinors}

For later use in lower-dimensional QFT, take
\[
  \eta_{\mu\nu}=\operatorname{diag}(-,+,+),
  \qquad
  \Gamma^0=\ii\sigma^2,
  \qquad
  \Gamma^1=\sigma^1,
  \qquad
  \Gamma^2=\sigma^3 .
\]
Then
\[
  \{\Gamma^\mu,\Gamma^\nu\}=2\eta^{\mu\nu},
\]
and all three matrices are real.  Therefore a Majorana spinor may be taken
to have real components.  Raise and lower two-component indices by
\[
  \chi^a=\epsilon^{ab}\chi_b,
  \qquad
  \epsilon^{12}=\epsilon_{12}=1 .
\]
The lowered gamma matrices
\[
  \Gamma^\mu_{ab}:=(\Gamma^\mu)_a{}^c\epsilon_{cb}
\]
are symmetric in \(a,b\).  For two Grassmann-even Majorana spinors
\(\zeta_1,\zeta_2\),
\[
  \bar\zeta_1\Gamma^\mu\zeta_2
  =
  \zeta_1^T\ii\Gamma^0\Gamma^\mu\zeta_2
  =
  \ii\,\zeta_1^a\Gamma^\mu_{ab}\zeta_2^b .
\]

\section{Two-Dimensional Lorentzian Spinors}

In two-dimensional Lorentzian calculations the monograph uses
\[
  \eta_{\mu\nu}=\operatorname{diag}(-,+),
  \qquad
  \{\gamma^\mu,\gamma^\nu\}=2\eta^{\mu\nu},
  \qquad
  \epsilon^{01}=+1 .
\]
The chirality matrix is
\[
  \gamma=\gamma^0\gamma^1,
  \qquad
  \gamma^2=1,
  \qquad
  \{\gamma,\gamma^\mu\}=0 .
\]
The trace convention compatible with the anomaly calculation is
\[
  \operatorname{tr}(\gamma^\mu\gamma\,\gamma^\nu)
  =
  2\epsilon^{\mu\nu}.
\]
As in four dimensions, \(\bar\psi=\ii\psi^\dagger\gamma^0\) in the
mostly-plus convention.

\section{Euclidean Spinors}

Euclidean gamma matrices in even dimension \(d\) are denoted
\(\Gamma^1,\ldots,\Gamma^d\) and obey
\[
  \{\Gamma^i,\Gamma^j\}=2\delta^{ij}.
\]
A concrete recursive construction starts in \(d=2\) with
\[
  \Gamma^1_{(2)}=\sigma_1,
  \qquad
  \Gamma^2_{(2)}=\sigma_2,
\]
and for even \(d>2\) sets
\[
  \Gamma^\mu_{(d)}
  =
  \Gamma^\mu_{(d-2)}\otimes(-\sigma_3),
  \qquad
  \mu=1,\ldots,d-2,
\]
\[
  \Gamma^{d-1}_{(d)}
  =
  I_{2^{d/2-1}}\otimes\sigma_1,
  \qquad
  \Gamma^d_{(d)}
  =
  I_{2^{d/2-1}}\otimes\sigma_2 .
\]
The Euclidean chirality matrix is
\[
  \Gamma_E
  =
  \ii^{-d/2}\Gamma^1\cdots\Gamma^d,
  \qquad
  \Gamma_E^2=1 .
\]
One convenient charge-conjugation matrix is
\[
  C_E=\Gamma^1\Gamma^3\cdots\Gamma^{d-1},
  \qquad d\equiv0\pmod4,
\]
and
\[
  C_E=\Gamma_E\Gamma^1\Gamma^3\cdots\Gamma^{d-1},
  \qquad d\equiv2\pmod4.
\]
These formulae fix a basis; invariant statements should be phrased in terms
of the corresponding real, complex, or quaternionic structure on the
Euclidean spinor module.

\section{Six- and Eight-Dimensional Euclidean Blocks}

For later supersymmetric and internal-symmetry applications, it is useful to
record block conventions for Euclidean \(SO(6)\) and \(SO(8)\) spinors.

In \(d=8\), the chiral and antichiral spinor representations are both real
eight-dimensional.  In a basis in which the charge-conjugation matrix is the
identity, write the off-diagonal gamma-matrix components as
\[
  (\Gamma^i)^{a\dot b}=\mathcal C^{i a\dot b},
  \qquad
  (\Gamma^i)^{\dot b a}=\mathcal C^{i\dot b a},
  \qquad
  \mathcal C^{i\dot b a}=-\mathcal C^{i a\dot b}.
\]
They obey
\[
  \mathcal C^{i a\dot c}\mathcal C^{j\dot c b}
  +
  \mathcal C^{j a\dot c}\mathcal C^{i\dot c b}
  =
  2\delta^{ij}\delta^{ab},
\]
\[
  \mathcal C^{i\dot a c}\mathcal C^{j c\dot b}
  +
  \mathcal C^{j\dot a c}\mathcal C^{i c\dot b}
  =
  2\delta^{ij}\delta^{\dot a\dot b},
\]
\[
  \mathcal C^{i a\dot c}\mathcal C^{i\dot d b}
  +
  \mathcal C^{i b\dot c}\mathcal C^{i\dot d a}
  =
  2\delta^{ab}\delta^{\dot c\dot d}.
\]
The first two identities are the Clifford relations; the third records the
triality symmetry that exchanges vector and spinor representations.

In \(d=6\), use \(\operatorname{Spin}(6)\simeq SU(4)\).  Chiral and
antichiral spinors are the fundamental and antifundamental representations
of \(SU(4)\).  Let \(I,J,K,L=1,\ldots,4\).  The gamma-matrix blocks are
antisymmetric tensors
\[
  \mathcal C^i_{IJ}=-\mathcal C^i_{JI},
  \qquad
  \mathcal C^{iIJ}
  =
  \frac12\epsilon^{IJKL}\mathcal C^i_{KL}
  =
  \left(\mathcal C^i_{JI}\right)^* .
\]
They satisfy
\[
  \sum_{i=1}^6
  \mathcal C^i_{IJ}\mathcal C^{iKL}
  =
  -2
  \left(\delta_I{}^K\delta_J{}^L-\delta_I{}^L\delta_J{}^K\right).
\]
These identities are the ones used when four-dimensional supersymmetry or
higher-dimensional reductions are translated into \(SU(4)\) index notation.
